\chapter{Layer 10: Automated Proof Completeness}
\label{ch:completeness}

\section{Purpose and Contract}
\label{sec:comp-purpose}

Layer~10 closes the loop. Its inbound contract is ``a set of obligations, some
still under \texttt{sorry}''; its outbound contract is ``an updated
completeness report.'' It is the system's conscience: it refuses to let open
goals be forgotten.

\section{The Sorry Database}
\label{sec:comp-sorrydb}

\texttt{sorrydb/} is the ledger of open goals. Each entry records an
obligation id, its source layer, and its status. Its \texttt{metadata.json}
carries \texttt{family\_id: 8} and a \texttt{corpora\_path} of
\texttt{formal\_research/snapkitty\_proofs/sorrydb/}, marking it as the
curated store of proof gaps. A repository's completeness score is:
\begin{equation}
C = 1 - \frac{\#\text{obligations under \texttt{sorry}}}{\#\text{total obligations}}.
\end{equation}

\section{The Sorry Hunter}
\label{sec:comp-hunter}

\texttt{layers/sorryhunter/} is the active closer. Its \texttt{README}
states:
\begin{quote}
\textbf{Layer 10: Automated Sorry Closing} \\
Purpose: Symbolic kernel + tactics + FFI + malice + axiom to close Lean~4
sorries. \\
Status: Active.
\end{quote}
The hunter combines four strategies:
\begin{enumerate}
  \item \textbf{Symbolic kernel} --- apply the verified C{--} primitives to
        compute a witness.
  \item \textbf{Tactics} --- invoke Lean's tactic language to close routine
        goals.
  \item \textbf{FFI} --- cross-check the kernel result against Lean
        (\cref{ch:ffi}) to certify the witness.
  \item \textbf{Malice + axiom} --- use the adversarial layer
        (\cref{sec:pnp-malice}) to stress-test and the axiom store to supply
        lemmas.
\end{enumerate}

\section{The Theorem Corpus}
\label{sec:comp-theorems}

The verified theorems to date (\texttt{theorems/THEOREMS.md}) are the closed
forms that anchor the calculator:
\begin{align}
\text{SumSquares} &: \sum_{k=1}^{n} k^{2} = \frac{n(n+1)(2n+1)}{6}, \\
\text{SumCubes}   &: \sum_{k=1}^{n} k^{3} = \left(\frac{n(n+1)}{2}\right)^{2}, \\
\text{SumLinear}  &: \sum_{k=1}^{n} k   = \frac{n(n+1)}{2}.
\end{align}
Each is both a Lean \texttt{theorem} and a C{--} closed-form entry, so the
sorry hunter can validate one against the other.

\section{Completeness as a First-Class Metric}
\label{sec:comp-metric}

By treating completeness as a measurable, reported quantity, \textsc{mathlib5}
converts the usual invisible debt of unproven lemmas into a visible,
monotonic-ish progress signal. The governance kernel's ``no component advances
unless coherent'' rule (\cref{ch:coherence}) can be specialized: a release is
permitted only if $C \ge C_{\min}$ for an agreed threshold. This binds the
abstract coherence model to a concrete proof metric.

\section{A Completeness Trace}
\label{sec:comp-trace}

We illustrate the completeness score $C$ on the current certified corpus.
Suppose the target theorem set is
\[
\mathcal{T} = \{\, T_1, T_2, T_3, T_4, T_5 \,\},
\]
where $T_1,T_2,T_3$ are the three closed forms currently proven in Lean
(\cref{ch:lean}) and $T_4,T_5$ are desired but still marked
\texttt{sorry} in \texttt{sorrydb}. Then
\[
C = \frac{|\{T_1,T_2,T_3\}|}{|\mathcal{T}|} = \frac{3}{5} = 0.6.
\]
The build gate (\cref{ch:build}) treats the \texttt{sorry} entries as open
items; the P/NP swarm (\cref{ch:pnp}) exposes $T_4,T_5$ as claimable
problems. When an agent submits a verified witness, $C$ rises to $0.8$,
then $1.0$. The metric is therefore a live dashboard of the system's
verification health, not a static claim.

\subsection{Interpretation}
A $C < 1$ is not a defect; it is honest accounting. Traditional libraries
often hide unproven lemmas behind unchecked imports. MATHLIB5 makes the
gap explicit and machine-countable, so a reviewer can state precisely how
much of the system rests on certified ground versus open proof debt.
